\input{preamble}

\FirstPage{1}
\LastPage{x}

\pdfcrop
%\draft

\begin{document}

\udc{XXX.XXX}
\msc{YYYYY, ZZZZZ, DDDDD}

\rutitle{Введите название Вашей статьи здесь}
\rutitleshort{Короткое название статьи, если необходимо (eсли не помещается сверху)}

\entitle{Enter your paper title here}
\entitleshort{Enter the short paper title here if necessary}

\ruauthor{А.~А.~Иванов\affil{1,2}, В.~В.~Смирнов\affil{2}}{И\,в\,а\,н\,о\,в~А.~В., С\,м\,и\,р\,н\,о\,в~В.~В.}
\enauthor{A.~A.~Ivanov\affil{1,2}, V.~V.~Smirnov\affil{2}}{I\,v\,a\,n\,o\,v~A.~V., S\,m\,i\,r\,n\,o\,v~V.~V.}

\ruaffil{Московский государственный университет имени М.~В.~Ломоносова,\\ 
Научно-исследовательский институт механики, \\
Россия, 119192, Москва, Мичуринский проспект, 1.
\and 
Самарский государственный технический университет,\\
Россия, 443100, Самара, ул. Молодогвардейская, 244.}

\enaffil{Lomonosov Moscow State University,\\ 
Institute of Mechanics, \\
1, Michurinsky prospekt, Moscow, 119192, Russian Federation.
\and 
Samara State Technical University,\\
244, Molodogvardeyskaya st., Samara, 443100, Russian Federation.}

\ruauthorsdetails{2}{% Количество авторов
\fullauthor{Андрей Алексеевич Иванов}
\CAuthor % Автор-корреспондент
\orcid{0000-0000-0000-xxxx} % ORCID ID автора 
\regalia{кандидат физико-математических наук, доцент} 
\position{старший научный сотрудник} 
\dept{лаб. упругости и пластичности} 
\email{anton@ya-mail.ru}
\fullauthor{Владимир Викторович Смирнов}
\orcid{0000-0000-0000-xxxx} % ORCID ID автора 
\regalia{доктор физико-математических наук, профессор} 
\position{заведующий кафедрой} 
\dept{каф. упругости и пластичности} 
\email{vlad_smirnov@ya-mail.ru}
}

\enauthorsdetails{2}{
\fullauthor{Andrew A. Ivanov}
\CAuthor % Сorresponding Author
\orcid{0000-0000-0000-xxxx} % ORCID ID
\regalia{Cand. Phys. \& Math. Sci., Associate Professor} 
\position{Senior Researcher} 
\dept{Lab. of Elasticity and Plasticity}
\email{anton@ya-mail.ru}
\fullauthor{Vladimir V. Smirnov}
\orcid{0000-0000-0000-xxxx} % ORCID ID avtora 
\regalia{Dr. Phys. \& Math. Sci., Professor} 
\position{Head of Department} 
\dept{Dept. of Elasticity and Plasticity}
\email[br]{vlad_smirnov@ya-mail.ru}
}

\rukeywords{ключевое слово, ключевая фраза, новая мысль}
\enkeywords{keyword, key phrase, new thought}

\ruabstract{Аннотация на русском языке должна содержать 3–5 предложений (не более 500 знаков без пробелов). Рекомендуется структурировать аннотацию: объект исследования, цель, использованные методы и подходы, основные результаты. Аннотация должна быть информативной, отражать основное содержание статьи и результаты исследований, следовать логике описания результатов в статье и не содержать общих слов. Аннотация не должна содержать ссылок на литературу и аббревиатуры (если это возможно).}

\enabstract{An abstract is a succinct summary of a longer piece of work, usually academic in nature, which is published in isolation from the main text and should therefore stand on its own and be understandable without reference to the longer piece. It should report the latter's essential facts, and should not exaggerate or contain material that is not there. Its purpose is to act as a reference tool (for example in a library abstracting service), enabling the reader to decide whether or not to read the full text. 

The abstract should reflect only what appears in the original paper. The Abstract should:
\begin{itemize}
\item Describe the main objective(s) of the study;
\item Explain how the study was done, including any model organisms used, without methodological detail;
\item Summarize the most important results and their significance.
\end{itemize}
Abstracts should not include citations and abbreviations, if possible.}

\makeentitle

\Section[n]{Introduction}
Here goes the introduction.

\Section{Section 1}
Example of formula: $x^2+y^2=z^2$.

\begin{figure}[h!]
   \centering
   \includegraphics[scale=.6]{Kelvin570g}
    \caption{The figure caption
    \label{fig:Kelvin570g}
    }
\end{figure}


\Section{Section 2}

Text text text text text. 
Text text text text text.

Text text text text text.

Text text text text text.
Text text text text text.


%example for  point list

\begin{itemize}
\item example for point list
\end {itemize}

Text text text text text. 
Text text text text text.

Text text text text text.

Text text text text text.
Text text text text text.


%example for numbered list
    \begin{enumerate}
    \item for numbered list
 
    \end{enumerate}

\Section[n]{Conclusion}

Text text text text text. 
Text text text text text.

Text text text text text.

Text text text text text.
Text text text text text.

\AdditionalContent{Acknowledgments}{This work was supported by the Russian Foundation for Basic Research (project no. 17–0x–xxxxx-a). }

\begin{thebibliography}{10}

\Bibitem{Kol83}
\by Kolotnikova~O.~V.
\paper Effectiveness of hardening by methods of plastic surface deformation of components operating at high temperatures
\jour Strength of Material
\vol 15
\issue 2
\yr 1983
\pages 292–295
\crossref{10.1007/BF01523487}


\Bibitem{ulg:bib:Birger1963}
\lang In Russian
\by Birger~I.~A.
\book  Ostatochnye napryazheniya \rm [Residual stresses]
\publ Mashgiz
\publaddr  Moscow
\yr  1963
\totalpages  262 

\Bibitem{ulg:bib:Uchaikin2013}
\by Uchaikin~V.~V.
\paper Heredity and Nonlocality
\inbook Fractional Derivatives for Physicists and Engineers
\yr 2013
\pages 3--58
\publ Springer Berlin Heidelberg
\publaddr Berlin
\bookvol 1
\volinfo Background and Theory
\crossref{10.1007/978-3-642-33911-0_1}



\end{thebibliography}

\newpage

\makerutitle

\end{document}
